\section{术语消化与总结}

本讲把表示、连续生成目标、Transformer 改造、数据 curriculum 和多任务应用压在同一系统里。若只记模型名，很容易混淆哪个组件解决哪个瓶颈。先用术语表重新按“问题--机制--证据”组织，再给出端到端检查表。

\subsection{高密度术语表}

\begin{center}
\small
\begin{tabularx}{\textwidth}{p{0.18\textwidth} X X}
\toprule
\textbf{术语} & \textbf{解决的问题} & \textbf{核心机制与课程关系} \\
\midrule
TAE & 像素视频序列太长 & 时间、高度、宽度各压缩 8 倍，用连续 latent 统一图像与视频。 \\
Latent & 哪个随机变量交给生成器 & 保留生成相关信息并舍弃部分像素冗余；压缩率决定质量--成本折中。 \\
Flow Matching & 如何从噪声生成媒体 & 学习连续时间 velocity field，推理用 ODE solver 积分。 \\
Cross-attention & 视频 token 如何读取文本 & 视频作 query，三个文本表示作 key/value，文本保持干净条件。 \\
AdaLN & timestep 如何进入每层 & 时间 embedding 生成 LayerNorm scale/shift，调制同一 block。 \\
MHA & 双向时空交互 & 每个 query head 有独立 K/V heads，适合非自回归全序列更新。 \\
GQA & 自回归 KV cache 太大 & 多个 query heads 共享 K/V；Movie Gen T2V 不依赖该优势。 \\
Context parallelism & 73K tokens 单卡放不下 & 按 token/context 维度分片并用集合通信交换注意力信息。 \\
Curriculum & 直接高分辨率训练太贵 & 从图像和低分辨率视频逐步提高分辨率、时长与任务复杂度。 \\
Net Win Rate & 自动指标不足 & 人类成对偏好中 A 胜比例减 B 胜比例，并报告 bootstrap 区间。 \\
Post-training & pretraining 覆盖不等于可控 & 用高质量任务数据提升运动、美学、编辑或身份保持。 \\
World model & “看起来合理”到底有多强 & 需区分视觉 plausibility、反事实一致性与可规划因果模型。 \\
\bottomrule
\end{tabularx}
\end{center}

\subsection{端到端工程检查表}

\begin{enumerate}
  \item \textbf{Representation：}时空压缩率、latent channels 与 patch size 是否完整报告？
  \item \textbf{Objective：}预测 noise、velocity 还是 data；solver 与采样步数是什么？
  \item \textbf{Conditioning：}文本 encoder、cross-attention、AdaLN 和 reference conditions 如何分工？
  \item \textbf{Context：}分辨率、帧率、时长对应多少 tokens，attention 与激活成本多大？
  \item \textbf{Data：}视觉、运动、内容、去重与 captioning 过滤是否可复现？
  \item \textbf{Curriculum：}哪些阶段训练图像、低分辨率视频、高分辨率视频和任务分支？
  \item \textbf{Evaluation：}是否同时报告失败样本、固定 prompt、人类偏好、运动与文本遵循？
  \item \textbf{Deployment：}权重、ODE steps、多卡通信、峰值显存和 serving 证据是否公开？
\end{enumerate}

\begin{importantbox}{全讲核心结论}
Movie Gen 的贡献可以压缩为一条系统链：用 TAE 把媒体压缩为 token；用 Flow Matching 定义连续运输目标；用随机初始化的 Llama-style 双向 Transformer 承担规模化 backbone；用数据清洗与渐进 curriculum 把 73K-token 训练变得可行；再用独立后训练扩展编辑、个性化和音频。能力提升真实存在，但复杂交互、长视频、评估、reasoning verifier 与 serving 仍是开放边界。
\end{importantbox}

